\ifdefined\inMainDoc\else
\documentclass[12pt,a4paper]{article}
% ================================================================
%  PRÉAMBULE COMMUN - ANNALES CORRIGÉES
%  Université de Kara - FaST - Licence Fondamentale MPC
%  Design optimisé impression noir et blanc
% ================================================================

% -- ENCODAGE & LANGUE --
\usepackage[utf8]{inputenc}
\usepackage[T1]{fontenc}
\usepackage[french]{babel}
\usepackage{lmodern}
\usepackage{microtype}

% -- MISE EN PAGE --
\usepackage[a4paper, top=2.8cm, bottom=2.5cm, left=2.5cm, right=2.5cm, headheight=14pt]{geometry}
\usepackage{setspace}
\setstretch{1.2}
\usepackage{parskip}
\setlength{\parindent}{1.2em}
\setlength{\parskip}{4pt}

% -- MATHS --
\usepackage{amsmath, amssymb, mathtools, amsthm}
\usepackage{mathrsfs}
\usepackage{dsfont}
\usepackage{bm}
\usepackage{cancel}
\usepackage{textcomp}
% Définition de \oiint (intégrale de surface fermée) sans package supplémentaire
\newcommand{\oiint}{\mathop{\iint\!\!\!\!\!\!\!\bigcirc}\limits}

% -- PHYSIQUE & CHIMIE --
\usepackage{physics}
\usepackage{siunitx}
% Lorsque physics est chargé, siunitx omet \qty. On le restaure ici :
\AtBeginDocument{\RenewCommandCopy\qty\SI}
\sisetup{locale = FR, detect-all, output-decimal-marker={,}}
% Commande de substitution pour les unités posant problème avec pdflatex
% (ohm, degreeCelsius, micro, angstrom, percent utilisent des caractères non-ASCII)
\newcommand{\qtyohm}[1]{\ensuremath{\num{#1}\,\Omega}}
\newcommand{\qtydegc}[1]{\ensuremath{\num{#1}\,{}^\circ\text{C}}}
\newcommand{\qtyangstrom}[1]{\ensuremath{\num{#1}\,\text{\AA}}}
\newcommand{\qtypercent}[1]{\ensuremath{\num{#1}\,\%}}
\newcommand{\qtyum}[1]{\ensuremath{\num{#1}\,\mu\text{m}}}
\usepackage[version=4]{mhchem}

% -- GRAPHISMES --
\usepackage{graphicx}
\usepackage{float}
\usepackage{caption}
\usepackage{subcaption}
\usepackage{wrapfig}
\usepackage{tikz}
\usetikzlibrary{calc, arrows.meta, patterns, shapes.geometric}
\usepackage{pgfplots}
\pgfplotsset{compat=1.18}

% -- TABLEAUX --
\usepackage{booktabs}
\usepackage{array}
\usepackage{tabularx}
\usepackage{multirow}

% -- LISTES --
\usepackage{enumitem}
\setlist[enumerate]{topsep=3pt, itemsep=2pt, parsep=0pt}
\setlist[itemize]{topsep=3pt, itemsep=2pt, parsep=0pt, label={.}}

% -- BOÎTES (noir et blanc) --
\usepackage{mdframed}
\usepackage{tcolorbox}
\tcbuselibrary{skins, breakable, theorems}

% Boîte EXERCICE : encadré avec filet épais, fond blanc
\newmdenv[
  linewidth=1.2pt,
  linecolor=black,
  roundcorner=3pt,
  skipabove=10pt,
  skipbelow=6pt,
  innertopmargin=8pt,
  innerbottommargin=8pt,
  innerleftmargin=10pt,
  innerrightmargin=10pt,
  backgroundcolor=white,
  frametitle={\bfseries\small Exercice},
  frametitlealignment=\raggedright,
  frametitleaboveskip=4pt,
  frametitlebelowskip=4pt,
]{boiteexercice}

% Boîte CORRIGÉ : fond gris très clair + filet fin
\newmdenv[
  linewidth=0.6pt,
  linecolor=black,
  leftline=true,
  rightline=false,
  topline=false,
  bottomline=false,
  linecolor=black,
  innerleftmargin=12pt,
  innerrightmargin=8pt,
  innertopmargin=6pt,
  innerbottommargin=6pt,
  backgroundcolor=gray!8,
  skipabove=4pt,
  skipbelow=8pt,
]{boitecorrige}

% Boîte RÉSUMÉ DE COURS : double encadré
\newmdenv[
  linewidth=0.8pt,
  linecolor=black,
  roundcorner=2pt,
  backgroundcolor=gray!12,
  innertopmargin=10pt,
  innerbottommargin=10pt,
  innerleftmargin=12pt,
  innerrightmargin=12pt,
  skipabove=12pt,
  skipbelow=8pt,
]{boiteresume}

% Boîte ATTENTION/REMARQUE : barre gauche épaisse
\newmdenv[
  leftline=true,
  rightline=false,
  topline=false,
  bottomline=false,
  linewidth=3pt,
  linecolor=black,
  backgroundcolor=gray!5,
  innerleftmargin=12pt,
  innerrightmargin=8pt,
  innertopmargin=5pt,
  innerbottommargin=5pt,
  skipabove=6pt,
  skipbelow=6pt,
]{boiteremarque}

% -- DIFFICULTÉ DES EXERCICES --
\newcommand{\facile}{$\star$}
\newcommand{\moyen}{$\star\!\star$}
\newcommand{\difficile}{$\star\!\star\!\star$}
\newcommand{\niveau}[1]{\hfill\textbf{#1}\hspace{2pt}}

% -- EN-TÊTES ET PIEDS DE PAGE --
\usepackage{fancyhdr}
\pagestyle{fancy}
\fancyhf{}
\renewcommand{\headrulewidth}{0.5pt}
\renewcommand{\footrulewidth}{0.3pt}

% -- HYPERLIENS --
\usepackage[hidelinks]{hyperref}
\hypersetup{
    colorlinks=false,
    pdftitle={Annale Corrigée - Licence MPC - Université de Kara}
}

% -- TITRES --
\usepackage{titlesec}
\titleformat{\section}{\Large\bfseries}{\thesection.}{0.8em}{}[\vspace{2pt}\hrule\vspace{4pt}]
\titleformat{\subsection}{\large\bfseries}{\thesubsection.}{0.7em}{}
\titleformat{\subsubsection}{\normalsize\bfseries}{\thesubsubsection.}{0.6em}{}

% -- THÉORÈMES --
\theoremstyle{definition}
\newtheorem{defn}{Définition}[section]
\newtheorem{prop}{Propriété}[section]
\newtheorem{thm}{Théorème}[section]
\newtheorem{corol}{Corollaire}[section]
\newtheorem{exemple}{Exemple}[section]

% -- COMMANDES MATHÉMATIQUES UTILES --
\newcommand{\vect}[1]{\overrightarrow{#1}}
\newcommand{\R}{\mathbb{R}}
\newcommand{\N}{\mathbb{N}}
\newcommand{\Z}{\mathbb{Z}}
\newcommand{\C}{\mathbb{C}}
\renewcommand{\d}{\mathrm{d}}
\newcommand{\deriv}[2]{\dfrac{\d #1}{\d #2}}
\newcommand{\dpart}[2]{\dfrac{\partial #1}{\partial #2}}
\newcommand{\rot}{\overrightarrow{\mathrm{rot}}\,}
\renewcommand{\div}{\mathrm{div}\,}
\renewcommand{\grad}{\overrightarrow{\mathrm{grad}}\,}
\newcommand{\e}{\mathrm{e}}
\newcommand{\ii}{\mathrm{i}}
\renewcommand{\Tr}{\mathrm{Tr}}

% Numérotation des équations par section
\numberwithin{equation}{section}

% -- RÉSULTATS ENCADRÉS --
\newcommand{\res}[1]{\[\boxed{#1}\]}

% -- REMARQUE INLINE --
\newcommand{\rmq}[1]{\begin{boiteremarque}\textbf{Remarque.} #1\end{boiteremarque}}

% -- MISE EN VALEUR DU RÉSULTAT FINAL --
\newcommand{\reponse}[1]{%
  \begin{center}
    \fbox{\begin{minipage}{0.92\textwidth}\centering #1\end{minipage}}
  \end{center}%
}


\lhead{\small\textit{Algèbre Linéaire - Licence 1}}
\rhead{\small\textit{FaST, Université de Kara}}
\cfoot{\thepage}

\begin{document}
\fi

\begin{center}
  {\LARGE\bfseries ALGÈBRE LINÉAIRE}\\[0.3cm]
  {\normalsize Licence 1 - Mathématiques, Physique et Chimie}\\[0.1cm]
  \rule{0.7\textwidth}{0.8pt}
\end{center}

\section*{Résumé du cours}

\begin{boiteresume}

\textbf{1. Espaces vectoriels.} Un ensemble $E$ muni d'une addition et d'une multiplication scalaire est un espace vectoriel sur $\mathbb{R}$ s'il vérifie les 8 axiomes (associativité, distributivité, élément neutre $\vec{0}$, opposé, etc.). Une famille $(\vec{e}_1,\ldots,\vec{e}_n)$ est une \textit{base} de $E$ si elle est libre (aucun vecteur combinaison des autres) et génératrice ($E = \text{vect}(\vec{e}_1,\ldots,\vec{e}_n)$). La dimension de $E$ est le nombre de vecteurs dans toute base.

\textbf{2. Applications linéaires.} Une application $f : E \to F$ est linéaire si $f(\alpha\vec{u}+\beta\vec{v}) = \alpha f(\vec{u})+\beta f(\vec{v})$ pour tous $\alpha,\beta\in\mathbb{R}$ et $\vec{u},\vec{v}\in E$. Le noyau $\ker f = \{\vec{x}\in E \mid f(\vec{x})=\vec{0}\}$ et l'image $\text{Im}\,f = \{f(\vec{x}) \mid \vec{x}\in E\}$. Théorème du rang : $\dim(\ker f) + \dim(\text{Im}\,f) = \dim E$.

\textbf{3. Matrices.} Une matrice $A\in\mathcal{M}_{m,n}(\mathbb{R})$ représente une application linéaire dans des bases. Le produit $AB$ est défini si le nombre de colonnes de $A$ égale le nombre de lignes de $B$ : $(AB)_{ij} = \sum_k a_{ik}b_{kj}$. Une matrice $A$ est inversible si et seulement si $\det A \neq 0$.

\textbf{4. Déterminants.} Pour $A\in\mathcal{M}_n(\mathbb{R})$, le déterminant vérifie : $\det(AB)=\det A\cdot\det B$, $\det(A^T)=\det A$. Calcul par développement selon une ligne ou cofacteurs.

\textbf{5. Valeurs propres et vecteurs propres.} $\lambda$ est valeur propre de $A$ si $\exists\vec{x}\neq\vec{0}$ tel que $A\vec{x}=\lambda\vec{x}$. Les valeurs propres sont les racines du polynôme caractéristique $P_A(\lambda) = \det(A-\lambda I)$. La matrice est diagonalisable si et seulement si elle possède $n$ vecteurs propres linéairement indépendants.

\textbf{6. Systèmes linéaires.} Le système $A\vec{x}=\vec{b}$ est résolu par la méthode de Gauss (échelonnement). Il admet une solution unique si $\det A\neq 0$, une infinité de solutions si $A$ est singulière et $\vec{b}\in\text{Im}\,A$, ou aucune solution sinon.

\end{boiteresume}

\newpage
\section{Espaces vectoriels et bases}

\subsection*{Exercice 1 - Sous-espaces vectoriels \niveau{\facile}}

\begin{boiteexercice}
Dans $\mathbb{R}^3$, déterminer parmi les ensembles suivants lesquels sont des sous-espaces vectoriels :
\begin{enumerate}[label=(\alph*)]
  \item $F = \{(x,y,z)\in\mathbb{R}^3 \mid x + 2y - z = 0\}$
  \item $G = \{(x,y,z)\in\mathbb{R}^3 \mid x^2 + y^2 = 0\}$
  \item $H = \{(x,y,z)\in\mathbb{R}^3 \mid x + y + z = 1\}$
  \item $K = \{(x,y,z)\in\mathbb{R}^3 \mid x = y = 0\}$
\end{enumerate}
\end{boiteexercice}

\begin{boitecorrige}
Pour être un sous-espace vectoriel, un ensemble doit : (i) contenir $\vec{0}$, (ii) être stable par addition, (iii) être stable par multiplication scalaire.

\textbf{(a) $F$ est un sous-espace vectoriel.}
$\vec{0}$: $0+2(0)-0=0$. Stabilité: si $\vec{u},\vec{v}\in F$, alors $(x_1+2y_1-z_1)+(x_2+2y_2-z_2)=0+0=0$, donc $\vec{u}+\vec{v}\in F$. Scalaire: $\alpha(x_1+2y_1-z_1)=\alpha\cdot0=0$. C'est le plan de vecteur normal $(1,2,-1)$.

\textbf{(b) $G$ est un sous-espace vectoriel.}
$x^2+y^2=0$ avec $x,y\in\mathbb{R}$ implique $x=y=0$. Donc $G=\{(0,0,z)\mid z\in\mathbb{R}\}$, qui est l'axe $(Oz)$, bien un sous-espace.

\textbf{(c) $H$ n'est pas un sous-espace vectoriel.}
$\vec{0}\notin H$ car $0+0+0=0\neq1$.

\textbf{(d) $K$ est un sous-espace vectoriel.}
$K=\{(0,0,z)\mid z\in\mathbb{R}\}$ est l'axe $(Oz)$. Mêmes vérifications immédiates.
\end{boitecorrige}

\subsection*{Exercice 2 - Base et dimension \niveau{\facile}}

\begin{boiteexercice}
On considère les vecteurs $\vec{u}_1=(1,1,0)$, $\vec{u}_2=(0,1,1)$, $\vec{u}_3=(1,0,-1)$ de $\mathbb{R}^3$.
\begin{enumerate}
  \item Ces vecteurs forment-ils une base de $\mathbb{R}^3$ ?
  \item Exprimer $\vec{v}=(2,3,1)$ dans cette base.
\end{enumerate}
\end{boiteexercice}

\begin{boitecorrige}
\textbf{1.} On calcule le déterminant :
\[
\det\begin{pmatrix}1&0&1\\1&1&0\\0&1&-1\end{pmatrix}
= 1\cdot(1\cdot(-1)-0\cdot1) - 0 + 1\cdot(1\cdot1-1\cdot0)
= (-1) + (1) = -2 + 2
\]
Développons correctement : $1(1\cdot(-1)-0\cdot1) - 0(1\cdot(-1)-0\cdot0) + 1(1\cdot1-1\cdot0) = -1+0+1=0$.

Hmm, recalculons avec la matrice dont les colonnes sont $\vec{u}_1,\vec{u}_2,\vec{u}_3$ :
\[
\det\begin{pmatrix}1&0&1\\1&1&0\\0&1&-1\end{pmatrix}
= 1\begin{vmatrix}1&0\\1&-1\end{vmatrix} - 0 + 1\begin{vmatrix}1&1\\0&1\end{vmatrix}
= 1(-1-0) + 1(1-0) = -1+1 = 0
\]
Le déterminant est nul : les vecteurs ne forment \textbf{pas} une base de $\mathbb{R}^3$. Ils sont liés.

On vérifie : $\vec{u}_3 = \vec{u}_1 - \vec{u}_2$ (car $(1,0,-1)=(1,1,0)-(0,1,1)$). Les trois vecteurs sont donc coplanaires.

\textbf{2.} Puisque $(\vec{u}_1,\vec{u}_2,\vec{u}_3)$ n'est pas une base, on ne peut pas toujours exprimer $\vec{v}$ de façon unique. On cherche $\alpha,\beta,\gamma$ tels que $\alpha\vec{u}_1+\beta\vec{u}_2+\gamma\vec{u}_3=\vec{v}$ :
\[
\begin{cases}\alpha+\gamma=2\\\alpha+\beta=3\\\beta-\gamma=1\end{cases}
\]
De la 2e et 3e : $\beta=\gamma+1$, donc $\alpha = 3-\beta = 2-\gamma$. Substitution dans la 1re : $(2-\gamma)+\gamma=2$. La solution a un paramètre libre : $\gamma=t$, $\beta=t+1$, $\alpha=2-t$ pour tout $t\in\mathbb{R}$. Par exemple pour $t=0$ : $\vec{v} = 2\vec{u}_1 + \vec{u}_2 + 0\vec{u}_3$.
\end{boitecorrige}

\section{Matrices et déterminants}

\subsection*{Exercice 3 - Opérations matricielles \niveau{\facile}}

\begin{boiteexercice}
Soient $A = \begin{pmatrix}1&2\\3&4\end{pmatrix}$ et $B = \begin{pmatrix}-1&0\\2&1\end{pmatrix}$.
\begin{enumerate}
  \item Calculer $AB$, $BA$, $A+2B$ et $A^T$.
  \item Calculer $\det A$, $\det B$ et $\det(AB)$.
  \item La matrice $A$ est-elle inversible ? Si oui, calculer $A^{-1}$.
\end{enumerate}
\end{boiteexercice}

\begin{boitecorrige}
\textbf{1.}
\[
AB = \begin{pmatrix}1(-1)+2(2) & 1(0)+2(1)\\ 3(-1)+4(2) & 3(0)+4(1)\end{pmatrix}
= \begin{pmatrix}3&2\\5&4\end{pmatrix}
\]
\[
BA = \begin{pmatrix}(-1)(1)+0(3) & (-1)(2)+0(4)\\ 2(1)+1(3) & 2(2)+1(4)\end{pmatrix}
= \begin{pmatrix}-1&-2\\5&8\end{pmatrix}
\]
$AB \neq BA$ : le produit matriciel n'est pas commutatif en général.
\[
A+2B = \begin{pmatrix}1-2&2+0\\3+4&4+2\end{pmatrix} = \begin{pmatrix}-1&2\\7&6\end{pmatrix}, \quad A^T = \begin{pmatrix}1&3\\2&4\end{pmatrix}
\]

\textbf{2.} $\det A = 1\times4 - 2\times3 = 4-6 = -2$. $\det B = (-1)(1)-(0)(2) = -1$. $\det(AB) = \det A\cdot\det B = (-2)(-1) = 2$. Vérification : $\det\begin{pmatrix}3&2\\5&4\end{pmatrix}=12-10=2$.

\textbf{3.} $\det A = -2 \neq 0$, donc $A$ est inversible :
\[
A^{-1} = \frac{1}{-2}\begin{pmatrix}4&-2\\-3&1\end{pmatrix} = \begin{pmatrix}-2&1\\3/2&-1/2\end{pmatrix}
\]
Vérification : $AA^{-1} = \begin{pmatrix}1(- 2)+2(3/2) & 1(1)+2(-1/2)\\ 3(-2)+4(3/2) & 3(1)+4(-1/2)\end{pmatrix} = \begin{pmatrix}-2+3 & 1-1\\ -6+6 & 3-2\end{pmatrix} = I$.
\end{boitecorrige}

\subsection*{Exercice 4 - Résolution de systèmes linéaires \niveau{\moyen}}

\begin{boiteexercice}
Résoudre les systèmes suivants par la méthode de Gauss :
\begin{enumerate}[label=(\alph*)]
  \item $\begin{cases}2x + y - z = 3\\ x - y + 2z = -1\\ 3x + 2y + z = 4\end{cases}$

  \item $\begin{cases}x + 2y + z = 2\\ 2x + 4y + 2z = 4\\ x + y - z = 1\end{cases}$
\end{enumerate}
\end{boiteexercice}

\begin{boitecorrige}
\textbf{(a)} Matrice augmentée :
\[
\begin{pmatrix}2&1&-1&\mid&3\\1&-1&2&\mid&-1\\3&2&1&\mid&4\end{pmatrix}
\xrightarrow{L_1\leftrightarrow L_2}
\begin{pmatrix}1&-1&2&\mid&-1\\2&1&-1&\mid&3\\3&2&1&\mid&4\end{pmatrix}
\]
\[
\xrightarrow{L_2-2L_1,\;L_3-3L_1}
\begin{pmatrix}1&-1&2&\mid&-1\\0&3&-5&\mid&5\\0&5&-5&\mid&7\end{pmatrix}
\xrightarrow{L_3-\frac{5}{3}L_2}
\begin{pmatrix}1&-1&2&\mid&-1\\0&3&-5&\mid&5\\0&0&\frac{10}{3}&\mid&\frac{-4}{3}\end{pmatrix}
\]
De la 3e ligne : $z = \frac{-4/3}{10/3} = -\frac{2}{5}$. De la 2e : $3y = 5+5z = 5-2 = 3$, donc $y=1$. De la 1re : $x = -1+y-2z = -1+1+\frac{4}{5} = \frac{4}{5}$.
\[
\boxed{x = \frac{4}{5},\quad y = 1,\quad z = -\frac{2}{5}}
\]

\textbf{(b)} La 2e équation est $2\times$ la 1re (dépendante). Après élimination :
\[
\begin{pmatrix}1&2&1&\mid&2\\0&0&0&\mid&0\\0&-1&-2&\mid&-1\end{pmatrix}
\xrightarrow{L_2\leftrightarrow L_3}
\begin{pmatrix}1&2&1&\mid&2\\0&-1&-2&\mid&-1\\0&0&0&\mid&0\end{pmatrix}
\]
De la 2e : $y = 1-2z$. De la 1re : $x = 2-2y-z = 2-2(1-2z)-z = 3z$. Solution à un paramètre libre $z=t\in\mathbb{R}$ :
\[
\boxed{(x,y,z) = (3t,\;1-2t,\;t), \quad t\in\mathbb{R}}
\]
\end{boitecorrige}

\section{Valeurs propres et diagonalisation}

\subsection*{Exercice 5 - Valeurs propres \niveau{\moyen}}

\begin{boiteexercice}
Soit la matrice $A = \begin{pmatrix}3&1\\1&3\end{pmatrix}$.
\begin{enumerate}
  \item Calculer le polynôme caractéristique et les valeurs propres de $A$.
  \item Déterminer les vecteurs propres associés à chaque valeur propre.
  \item La matrice est-elle diagonalisable ? Si oui, écrire la décomposition $A = PDP^{-1}$.
\end{enumerate}
\end{boiteexercice}

\begin{boitecorrige}
\textbf{1.} $P_A(\lambda) = \det(A-\lambda I) = \det\begin{pmatrix}3-\lambda&1\\1&3-\lambda\end{pmatrix} = (3-\lambda)^2-1 = \lambda^2-6\lambda+8 = (\lambda-2)(\lambda-4)$.

Valeurs propres : $\lambda_1 = 2$ et $\lambda_2 = 4$.

\textbf{2.}
Pour $\lambda_1 = 2$ : $(A-2I)\vec{x}=\vec{0}$ donne $\begin{pmatrix}1&1\\1&1\end{pmatrix}\vec{x}=\vec{0}$, soit $x+y=0$. Vecteur propre : $\vec{v}_1 = \begin{pmatrix}1\\-1\end{pmatrix}$.

Pour $\lambda_2 = 4$ : $(A-4I)\vec{x}=\vec{0}$ donne $\begin{pmatrix}-1&1\\1&-1\end{pmatrix}\vec{x}=\vec{0}$, soit $-x+y=0$. Vecteur propre : $\vec{v}_2 = \begin{pmatrix}1\\1\end{pmatrix}$.

\textbf{3.} $A$ est symétrique réelle, donc diagonalisable (théorème spectral). Les deux vecteurs propres sont linéairement indépendants.
\[
P = \begin{pmatrix}1&1\\-1&1\end{pmatrix}, \quad D = \begin{pmatrix}2&0\\0&4\end{pmatrix}, \quad P^{-1} = \frac{1}{2}\begin{pmatrix}1&-1\\1&1\end{pmatrix}
\]
Vérification : $AP = PD$.
\end{boitecorrige}

\subsection*{Exercice 6 - Matrice $3\times3$ et diagonalisation \niveau{\moyen}}

\begin{boiteexercice}
Soit $A = \begin{pmatrix}2&1&0\\1&2&0\\0&0&3\end{pmatrix}$.
\begin{enumerate}
  \item Calculer les valeurs propres et les espaces propres.
  \item La matrice est-elle diagonalisable ?
  \item Calculer $A^n$ pour $n\in\mathbb{N}$.
\end{enumerate}
\end{boiteexercice}

\begin{boitecorrige}
\textbf{1.} $P_A(\lambda) = \det(A-\lambda I)$. En développant selon la 3e colonne :
\begin{align*}
P_A(\lambda) &= (3-\lambda)\det\begin{pmatrix}2-\lambda&1\\1&2-\lambda\end{pmatrix}\\
&= (3-\lambda)\bigl[(2-\lambda)^2-1\bigr]\\
&= (3-\lambda)(\lambda^2-4\lambda+3)\\
&= (3-\lambda)(3-\lambda)(1-\lambda) = (3-\lambda)^2(1-\lambda)
\end{align*}
Valeurs propres : $\lambda_1=1$ (simple) et $\lambda_2=3$ (double).

Pour $\lambda_1=1$ : $(A-I)\vec{x}=\vec{0}$ donne la matrice :
\[\begin{pmatrix}1&1&0\\1&1&0\\0&0&2\end{pmatrix}\]
soit $x+y=0$ et $z=0$.
\[
E_1 = \text{vect}\begin{pmatrix}1\\-1\\0\end{pmatrix}, \quad \dim E_1=1.
\]

Pour $\lambda_2=3$ : $(A-3I)\vec{x}=\vec{0}$ donne $\begin{pmatrix}-1&1&0\\1&-1&0\\0&0&0\end{pmatrix}$, soit $-x+y=0$.
\[
E_3 = \text{vect}\left\{\begin{pmatrix}1\\1\\0\end{pmatrix},\begin{pmatrix}0\\0\\1\end{pmatrix}\right\}, \quad \dim E_3=2.
\]

\textbf{2.} $\dim E_1+\dim E_3 = 1+2=3=\dim\mathbb{R}^3$. La matrice est \textbf{diagonalisable}.

\textbf{3.} Avec $P = \begin{pmatrix}1&1&0\\-1&1&0\\0&0&1\end{pmatrix}$, on a $A = PDP^{-1}$ avec $D=\text{diag}(1,3,3)$. Donc $A^n = PD^nP^{-1}$ avec $D^n=\text{diag}(1,3^n,3^n)$. Après calcul :
\[
A^n = \frac{1}{2}\begin{pmatrix}1+3^n&3^n-1&0\\3^n-1&1+3^n&0\\0&0&2\cdot3^n\end{pmatrix}
\]
\end{boitecorrige}

\subsection*{Exercice 7 - Application linéaire \niveau{\difficile}}

\begin{boiteexercice}
Soit $f : \mathbb{R}^3 \to \mathbb{R}^3$ l'application définie par $f(x,y,z) = (x+y, y+z, x+z)$.
\begin{enumerate}
  \item Vérifier que $f$ est linéaire.
  \item Déterminer $\ker f$ et $\text{Im}\,f$.
  \item Écrire la matrice de $f$ dans la base canonique.
  \item $f$ est-elle un isomorphisme ?
\end{enumerate}
\end{boiteexercice}

\begin{boitecorrige}
\textbf{1.} Soient $\vec{u}=(x_1,y_1,z_1)$, $\vec{v}=(x_2,y_2,z_2)$ et $\alpha,\beta\in\mathbb{R}$ :
\[
f(\alpha\vec{u}+\beta\vec{v}) = (\alpha x_1+\beta x_2+\alpha y_1+\beta y_2,\;\ldots)
= \alpha f(\vec{u})+\beta f(\vec{v})
\]
Les composantes sont des combinaisons linéaires des coordonnées, donc $f$ est bien linéaire.

\textbf{2.} $\ker f$ : $f(x,y,z)=\vec{0}$ donne
$\begin{cases}x+y=0\\y+z=0\\x+z=0\end{cases}$. De la 1re et 3e : $y=-x$, $z=-x$. Substitution dans la 2e : $-x+(-x)=-2x=0$, donc $x=0$. Ainsi $\ker f=\{\vec{0}\}$.

$\text{Im}\,f$ : par le théorème du rang, $\dim(\text{Im}\,f)=3-\dim(\ker f)=3-0=3$. Donc $\text{Im}\,f=\mathbb{R}^3$.

\textbf{3.} La matrice dans la base canonique est obtenue en appliquant $f$ aux vecteurs de base :
\[
[f] = \begin{pmatrix}1&1&0\\0&1&1\\1&0&1\end{pmatrix}
\]

\textbf{4.} $\ker f=\{\vec{0}\}$, donc $f$ est injective. Comme $\dim\mathbb{R}^3<\infty$ et $f$ est injective de $\mathbb{R}^3$ dans $\mathbb{R}^3$, elle est surjective, donc c'est un \textbf{isomorphisme}. On peut vérifier : $\det[f]=1(1-0)-1(0-1)+0=1+1=2\neq0$.
\end{boitecorrige}

\subsection*{Exercice 8 - Rang d'une matrice \niveau{\moyen}}

\begin{boiteexercice}
Calculer le rang des matrices suivantes :
\begin{enumerate}[label=(\alph*)]
  \item $A = \begin{pmatrix}1&2&3\\2&4&6\\1&1&2\end{pmatrix}$
  \item $B = \begin{pmatrix}1&0&2\\0&1&-1\\2&-1&5\end{pmatrix}$
  \item $C = \begin{pmatrix}1&2&0&1\\2&4&1&3\\0&0&1&1\end{pmatrix}$
\end{enumerate}
\end{boiteexercice}

\begin{boitecorrige}
\textbf{(a)} $L_2 \leftarrow L_2 - 2L_1$, $L_3 \leftarrow L_3 - L_1$ :
\[
\begin{pmatrix}1&2&3\\0&0&0\\0&-1&-1\end{pmatrix}
\xrightarrow{L_2\leftrightarrow L_3}
\begin{pmatrix}1&2&3\\0&-1&-1\\0&0&0\end{pmatrix}
\]
Deux lignes non nulles : $\text{rg}(A) = 2$.

\textbf{(b)} $L_3 \leftarrow L_3 - 2L_1$ puis $L_3 \leftarrow L_3 + L_2$ :
\[
\begin{pmatrix}1&0&2\\0&1&-1\\0&-1&1\end{pmatrix}
\to
\begin{pmatrix}1&0&2\\0&1&-1\\0&0&0\end{pmatrix}
\]
$\text{rg}(B) = 2$.

\textbf{(c)} On travaille sur la matrice $C$ (3 lignes, 4 colonnes). $L_2 \leftarrow L_2 - 2L_1$ :
\[
\begin{pmatrix}1&2&0&1\\0&0&1&1\\0&0&1&1\end{pmatrix}
\xrightarrow{L_3 - L_2}
\begin{pmatrix}1&2&0&1\\0&0&1&1\\0&0&0&0\end{pmatrix}
\]
$\text{rg}(C) = 2$.
\end{boitecorrige}

\subsection*{Exercice 9 - Inverse par la méthode de Gauss-Jordan \niveau{\moyen}}

\begin{boiteexercice}
Calculer l'inverse de la matrice $A = \begin{pmatrix}1&1&0\\0&1&1\\1&0&1\end{pmatrix}$ par la méthode de Gauss-Jordan.
\end{boiteexercice}

\begin{boitecorrige}
On augmente la matrice avec l'identité et on réduit $A$ à $I$ par opérations élémentaires :
\[
\left(\begin{array}{ccc|ccc}1&1&0&1&0&0\\0&1&1&0&1&0\\1&0&1&0&0&1\end{array}\right)
\xrightarrow{L_3-L_1}
\left(\begin{array}{ccc|ccc}1&1&0&1&0&0\\0&1&1&0&1&0\\0&-1&1&-1&0&1\end{array}\right)
\]
\[
\xrightarrow{L_3+L_2}
\left(\begin{array}{ccc|ccc}1&1&0&1&0&0\\0&1&1&0&1&0\\0&0&2&-1&1&1\end{array}\right)
\xrightarrow{L_3/2}
\left(\begin{array}{ccc|ccc}1&1&0&1&0&0\\0&1&1&0&1&0\\0&0&1&-1/2&1/2&1/2\end{array}\right)
\]
\[
\xrightarrow{L_2-L_3}
\left(\begin{array}{ccc|ccc}1&1&0&1&0&0\\0&1&0&1/2&1/2&-1/2\\0&0&1&-1/2&1/2&1/2\end{array}\right)
\xrightarrow{L_1-L_2}
\left(\begin{array}{ccc|ccc}1&0&0&1/2&-1/2&1/2\\0&1&0&1/2&1/2&-1/2\\0&0&1&-1/2&1/2&1/2\end{array}\right)
\]
\[
A^{-1} = \frac{1}{2}\begin{pmatrix}1&-1&1\\1&1&-1\\-1&1&1\end{pmatrix}
\]
Vérification : $\det A = 1(1-0) - 1(0-1) + 0 = 2 \neq 0$.
\end{boitecorrige}

\section{Systèmes linéaires avancés}

\subsection*{Exercice 10 - Règle de Cramer \niveau{\moyen}}

\begin{boiteexercice}
Résoudre le système suivant par la règle de Cramer :
\[
\begin{cases}2x + y - z = 5\\x - y + z = 0\\x + 2y + z = 3\end{cases}
\]
\end{boiteexercice}

\begin{boitecorrige}
La matrice du système est $A = \begin{pmatrix}2&1&-1\\1&-1&1\\1&2&1\end{pmatrix}$.

$\det A = 2(-1-2) - 1(1-1) + (-1)(2+1) = 2(-3) - 0 - 3 = -9$.

Matrices dérivées (on remplace chaque colonne par le second membre $\vec{b} = (5,0,3)^T$) :
\[
D_x = \det\begin{pmatrix}5&1&-1\\0&-1&1\\3&2&1\end{pmatrix} = 5(-1-2)-1(0-3)+(-1)(0+3) = -15+3-3=-15
\]
\[
D_y = \det\begin{pmatrix}2&5&-1\\1&0&1\\1&3&1\end{pmatrix} = 2(0-3)-5(1-1)+(-1)(3-0) = -6-0-3=-9
\]
\[
D_z = \det\begin{pmatrix}2&1&5\\1&-1&0\\1&2&3\end{pmatrix} = 2(-3-0)-1(3-0)+5(2+1) = -6-3+15=6
\]
\[
\boxed{x = \frac{D_x}{\det A} = \frac{-15}{-9} = \frac{5}{3}, \quad y = \frac{D_y}{\det A} = \frac{-9}{-9} = 1, \quad z = \frac{D_z}{\det A} = \frac{6}{-9} = -\frac{2}{3}}
\]
Vérification : $2(5/3)+1-(-2/3) = 10/3+1+2/3 = 12/3+1=5$.
\end{boitecorrige}

\subsection*{Exercice 11 - Déterminant d'une matrice $4\times4$ \niveau{\difficile}}

\begin{boiteexercice}
Calculer le déterminant de la matrice
\[
A = \begin{pmatrix}1&0&2&-1\\3&0&0&5\\2&1&4&-3\\1&0&5&0\end{pmatrix}
\]
par développement selon la colonne qui contient le plus de zéros.
\end{boiteexercice}

\begin{boitecorrige}
On développe selon la colonne 2 (trois zéros) :
\[
\det A = 0\cdot C_{12} + 0\cdot C_{22} + 1\cdot C_{32} + 0\cdot C_{42}
\]
Il ne reste que le cofacteur $C_{32} = (-1)^{3+2}M_{32} = -M_{32}$, où $M_{32}$ est le mineur obtenu en supprimant la ligne 3 et la colonne 2 :
\[
M_{32} = \det\begin{pmatrix}1&2&-1\\3&0&5\\1&5&0\end{pmatrix}
\]
On développe selon la 2e colonne (ou ligne 2) :
\[
= 1(0-25) - 2(0-5) + (-1)(15-0) = -25+10-15 = -30
\]
\[
\det A = 1\times(-(-30)) = 1\times30 = 30
\]

\textbf{Vérification partielle :} Le déterminant d'une matrice triangulaire supérieure serait le produit des termes diagonaux. Ce n'est pas le cas ici, mais on peut vérifier avec le développement selon la ligne 1 comme contrôle.
\end{boitecorrige}

\subsection*{Exercice 12 - Valeurs propres complexes et réduction \niveau{\difficile}}

\begin{boiteexercice}
Soit $A = \begin{pmatrix}0&-1&0\\1&0&0\\0&0&2\end{pmatrix}$.
\begin{enumerate}
  \item Calculer le polynôme caractéristique de $A$.
  \item Déterminer les valeurs propres (dans $\mathbb{C}$) et les sous-espaces propres.
  \item La matrice est-elle diagonalisable sur $\mathbb{R}$ ? Sur $\mathbb{C}$ ?
\end{enumerate}
\end{boiteexercice}

\begin{boitecorrige}
\textbf{1.} En développant selon la 3e colonne (terme $2-\lambda$ en bas à droite) :
\[
P_A(\lambda) = \det(A-\lambda I) = (2-\lambda)\det\begin{pmatrix}-\lambda&-1\\1&-\lambda\end{pmatrix}
= (2-\lambda)(\lambda^2+1)
\]

\textbf{2.} Valeurs propres : $\lambda_1 = 2$ (réelle) et $\lambda_{2,3} = \pm i$ (complexes conjuguées).

Pour $\lambda_1 = 2$ : $(A-2I)\vec{x}=\vec{0}$ : $\begin{pmatrix}-2&-1&0\\1&-2&0\\0&0&0\end{pmatrix}$, soit $-2x-y=0$ et $x-2y=0$. Ces deux équations donnent $x=y=0$, $z$ libre. $E_2 = \text{vect}\begin{pmatrix}0\\0\\1\end{pmatrix}$.

Pour $\lambda_2 = i$ (sur $\mathbb{C}$) : $(A-iI)\vec{x}=\vec{0}$ donne $-ix-y=0$, soit $y=-ix$. Sous-espace propre : $E_i = \text{vect}\begin{pmatrix}1\\-i\\0\end{pmatrix}$.

De même, $E_{-i} = \text{vect}\begin{pmatrix}1\\i\\0\end{pmatrix}$.

\textbf{3.} Les valeurs propres $\pm i$ ne sont pas réelles : $A$ \textbf{n'est pas diagonalisable sur $\mathbb{R}$}. En revanche, elle possède 3 valeurs propres complexes distinctes, donc elle est \textbf{diagonalisable sur $\mathbb{C}$}, avec $D = \text{diag}(2, i, -i)$.
\end{boitecorrige}

\subsection*{Exercice 13 - Espace vectoriel de polynômes \niveau{\moyen}}

\begin{boiteexercice}
Soit $E = \mathbb{R}_2[X]$ l'espace des polynômes de degré au plus 2, et $f : E \to \mathbb{R}^3$ définie par $f(P) = (P(0), P(1), P(2))$.
\begin{enumerate}
  \item Vérifier que $f$ est linéaire.
  \item Calculer $\ker f$ et $\text{Im}\, f$.
  \item En déduire si $f$ est un isomorphisme.
\end{enumerate}
\end{boiteexercice}

\begin{boitecorrige}
\textbf{1.} Pour $P, Q \in E$ et $\alpha, \beta \in \mathbb{R}$ :
\[
f(\alpha P + \beta Q) = ((\alpha P+\beta Q)(0), (\alpha P+\beta Q)(1), (\alpha P+\beta Q)(2))
\]
\[
= (\alpha P(0)+\beta Q(0), \alpha P(1)+\beta Q(1), \alpha P(2)+\beta Q(2)) = \alpha f(P) + \beta f(Q)
\]
Donc $f$ est linéaire.

\textbf{2.} $\ker f = \{P\in E : P(0)=P(1)=P(2)=0\}$. Un polynôme de degré $\leq 2$ ayant 3 racines distinctes est le polynôme nul. Donc $\ker f = \{0\}$.

Par le théorème du rang : $\dim(\text{Im}\,f) = \dim E - \dim(\ker f) = 3 - 0 = 3$. Donc $\text{Im}\,f = \mathbb{R}^3$ (surjective).

\textbf{3.} $f$ est injective ($\ker f = \{0\}$) et surjective. C'est donc un \textbf{isomorphisme}.

Cela signifie que tout triplet $(a, b, c) \in \mathbb{R}^3$ est l'image d'un unique polynôme de degré $\leq 2$ : c'est l'interpolation de Lagrange en $x=0,1,2$.
\end{boitecorrige}

\subsection*{Exercice 14 - Diagonalisation et puissance de matrice \niveau{\difficile}}

\begin{boiteexercice}
Soit $A = \begin{pmatrix}4&-2\\1&1\end{pmatrix}$.
\begin{enumerate}
  \item Calculer les valeurs propres et les vecteurs propres de $A$.
  \item Diagonaliser $A$ et exprimer $A^n$ pour tout $n\in\mathbb{N}$.
  \item Calculer $A^{10}$.
\end{enumerate}
\end{boiteexercice}

\begin{boitecorrige}
\textbf{1.} $P_A(\lambda) = (4-\lambda)(1-\lambda)+2 = \lambda^2-5\lambda+6 = (\lambda-2)(\lambda-3)$.

Valeurs propres : $\lambda_1 = 2$, $\lambda_2 = 3$.

Pour $\lambda_1 = 2$ : $(A-2I)\vec{x}=\vec{0}$ : $2x-2y=0$, soit $\vec{v}_1 = \begin{pmatrix}1\\1\end{pmatrix}$.

Pour $\lambda_2 = 3$ : $(A-3I)\vec{x}=\vec{0}$ : $x-2y=0$, soit $\vec{v}_2 = \begin{pmatrix}2\\1\end{pmatrix}$.

\textbf{2.} $P = \begin{pmatrix}1&2\\1&1\end{pmatrix}$, $D = \begin{pmatrix}2&0\\0&3\end{pmatrix}$, $P^{-1} = \frac{1}{\det P}\begin{pmatrix}1&-2\\-1&1\end{pmatrix} = \begin{pmatrix}-1&2\\1&-1\end{pmatrix}$.

$A^n = PD^nP^{-1} = \begin{pmatrix}1&2\\1&1\end{pmatrix}\begin{pmatrix}2^n&0\\0&3^n\end{pmatrix}\begin{pmatrix}-1&2\\1&-1\end{pmatrix}$

\[
= \begin{pmatrix}2^n&2\cdot3^n\\2^n&3^n\end{pmatrix}\begin{pmatrix}-1&2\\1&-1\end{pmatrix}
= \begin{pmatrix}-2^n+2\cdot3^n & 2^{n+1}-2\cdot3^n \\ -2^n+3^n & 2^{n+1}-3^n\end{pmatrix}
\]

Simplifions : $A^n = \begin{pmatrix}2\cdot3^n-2^n & 2^{n+1}-2\cdot3^n\\3^n-2^n & 2^{n+1}-3^n\end{pmatrix}$.

\textbf{3.} Pour $n=10$ : $2^{10} = 1024$, $3^{10} = 59049$ :
\[
A^{10} = \begin{pmatrix}2\times59049-1024 & 2048-2\times59049\\59049-1024 & 2048-59049\end{pmatrix}
= \begin{pmatrix}117074 & -116050\\58025 & -57001\end{pmatrix}
\]
\end{boitecorrige}


% ================================================================
\subsection*{Exercice 15 - Diagonalisation d'une matrice symétrique \niveau{\moyen}}
% ================================================================
\begin{boiteexercice}
Soit la matrice $A = \begin{pmatrix}5 & 2 \\ 2 & 5\end{pmatrix}$.

\begin{enumerate}
  \item Calculer le polynôme caractéristique de $A$ et en déduire ses valeurs propres.
  \item Déterminer les vecteurs propres associés à chaque valeur propre.
  \item Former la matrice de passage $P$ et la matrice diagonale $D = P^{-1}AP$.
  \item Vérifier que $A = PDP^{-1}$.
\end{enumerate}
\end{boiteexercice}

\begin{boitecorrige}
\textbf{1. Polynôme caractéristique :}
\[
\chi_A(\lambda) = \det(A - \lambda I) = \begin{vmatrix}5-\lambda & 2 \\ 2 & 5-\lambda\end{vmatrix}
= (5-\lambda)^2 - 4 = \lambda^2 - 10\lambda + 21 = (\lambda-3)(\lambda-7)
\]
Les valeurs propres sont $\lambda_1 = 3$ et $\lambda_2 = 7$.

\textbf{2. Vecteurs propres :}

\textit{Pour $\lambda_1 = 3$ :} $(A-3I)\vec{x} = \vec{0}$ :
\[
\begin{pmatrix}2 & 2 \\ 2 & 2\end{pmatrix}\begin{pmatrix}x\\y\end{pmatrix} = \vec{0} \;\Rightarrow\; x+y=0 \;\Rightarrow\; \vec{v}_1 = \begin{pmatrix}1\\-1\end{pmatrix}
\]

\textit{Pour $\lambda_2 = 7$ :} $(A-7I)\vec{x} = \vec{0}$ :
\[
\begin{pmatrix}-2 & 2 \\ 2 & -2\end{pmatrix}\begin{pmatrix}x\\y\end{pmatrix} = \vec{0} \;\Rightarrow\; -x+y=0 \;\Rightarrow\; \vec{v}_2 = \begin{pmatrix}1\\1\end{pmatrix}
\]

\textbf{3. Matrice de passage et diagonalisation :}
\[
P = \begin{pmatrix}1 & 1 \\ -1 & 1\end{pmatrix}, \qquad \det P = 1\cdot1-(-1)\cdot1 = 2
\]
\[
P^{-1} = \frac{1}{2}\begin{pmatrix}1 & -1 \\ 1 & 1\end{pmatrix}
\]
\[
D = P^{-1}AP = \begin{pmatrix}3 & 0 \\ 0 & 7\end{pmatrix}
\]

\textbf{4. Vérification :}
\[
PDP^{-1} = \begin{pmatrix}1&1\\-1&1\end{pmatrix}\begin{pmatrix}3&0\\0&7\end{pmatrix}\frac{1}{2}\begin{pmatrix}1&-1\\1&1\end{pmatrix}
= \frac{1}{2}\begin{pmatrix}3+7 & -3+7\\-3+7 & 3+7\end{pmatrix} = \begin{pmatrix}5&2\\2&5\end{pmatrix} = A \;
\]

\textbf{Remarque :} $A$ est symétrique réelle, donc ses vecteurs propres sont bien orthogonaux : $\vec{v}_1\cdot\vec{v}_2 = 1\cdot1+(-1)\cdot1 = 0$.
\end{boitecorrige}

% ================================================================
\subsection*{Exercice 16 - Forme de Jordan (matrice non diagonalisable) \niveau{\difficile}}
% ================================================================
\begin{boiteexercice}
Soit la matrice $A = \begin{pmatrix}2 & 1 \\ 0 & 2\end{pmatrix}$.

\begin{enumerate}
  \item Calculer le polynôme caractéristique. Quelle est la valeur propre (sa multiplicité algébrique) ?
  \item Trouver le sous-espace propre associé. Quelle est la multiplicité géométrique ?
  \item La matrice est-elle diagonalisable ? Justifier.
  \item Déterminer un vecteur propre généralisé $\vec{w}$ tel que $(A-\lambda I)\vec{w} = \vec{v}$ où $\vec{v}$ est le vecteur propre.
  \item Écrire la forme de Jordan $J$ de $A$ et la matrice de passage $P$ telle que $A = PJP^{-1}$.
\end{enumerate}
\end{boiteexercice}

\begin{boitecorrige}
\textbf{1. Polynôme caractéristique :}
\[
\chi_A(\lambda) = (2-\lambda)^2
\]
Valeur propre unique : $\lambda = 2$ de multiplicité algébrique $m_a = 2$.

\textbf{2. Sous-espace propre :}
\[
A - 2I = \begin{pmatrix}0&1\\0&0\end{pmatrix}
\]
$(A-2I)\vec{x}=\vec{0}$ donne $y=0$, donc $\vec{v}_1 = \begin{pmatrix}1\\0\end{pmatrix}$.
Multiplicité géométrique : $m_g = 1$ (un seul vecteur propre indépendant).

\textbf{3.} Puisque $m_g = 1 < m_a = 2$, la matrice \textbf{n'est pas diagonalisable}.

\textbf{4. Vecteur propre généralisé :}

On cherche $\vec{w}$ tel que $(A-2I)\vec{w} = \vec{v}_1$ :
\[
\begin{pmatrix}0&1\\0&0\end{pmatrix}\begin{pmatrix}w_1\\w_2\end{pmatrix} = \begin{pmatrix}1\\0\end{pmatrix} \;\Rightarrow\; w_2 = 1, \quad w_1 \text{ libre}
\]
On choisit $w_1 = 0$, soit $\vec{w} = \begin{pmatrix}0\\1\end{pmatrix}$.

\textbf{5. Forme de Jordan :}

La base de Jordan est $\mathcal{B} = (\vec{v}_1, \vec{w})$, d'où :
\[
P = \begin{pmatrix}1 & 0 \\ 0 & 1\end{pmatrix} = I, \qquad J = \begin{pmatrix}2 & 1 \\ 0 & 2\end{pmatrix} = A
\]
Dans ce cas, $A$ est \emph{déjà} sous forme de Jordan. On a bien $J = P^{-1}AP$ avec $P=I$.

\textbf{Propriété utile :} $A^n = \begin{pmatrix}2^n & n\cdot2^{n-1}\\0&2^n\end{pmatrix}$. En effet, $(J)^n = \begin{pmatrix}\lambda^n & n\lambda^{n-1}\\0&\lambda^n\end{pmatrix}$ pour un bloc de Jordan $2\times2$.
\end{boitecorrige}

% ================================================================
\subsection*{Exercice 17 - Calcul de $A^n$ par diagonalisation \niveau{\moyen}}
% ================================================================
\begin{boiteexercice}
Soit $A = \begin{pmatrix}2&1\\1&2\end{pmatrix}$.

\begin{enumerate}
  \item Calculer les valeurs propres $\lambda_1$ et $\lambda_2$ de $A$.
  \item Trouver les vecteurs propres associés.
  \item Former la matrice de passage $P$ et calculer $P^{-1}$.
  \item Exprimer $A^n = PD^nP^{-1}$ et simplifier.
  \item Application : calculer $A^5$.
\end{enumerate}
\end{boiteexercice}

\begin{boitecorrige}
\textbf{1. Valeurs propres :}
\[
\chi_A(\lambda) = (2-\lambda)^2 - 1 = \lambda^2 - 4\lambda + 3 = (\lambda-1)(\lambda-3)
\]
$\lambda_1 = 1$, $\lambda_2 = 3$.

\textbf{2. Vecteurs propres :}

\textit{Pour $\lambda_1 = 1$ :} $(A-I)\vec{x}=\vec{0}$ : $x+y=0$, donc $\vec{v}_1 = \begin{pmatrix}1\\-1\end{pmatrix}$.

\textit{Pour $\lambda_2 = 3$ :} $(A-3I)\vec{x}=\vec{0}$ : $-x+y=0$, donc $\vec{v}_2 = \begin{pmatrix}1\\1\end{pmatrix}$.

\textbf{3. Matrice de passage :}
\[
P = \begin{pmatrix}1&1\\-1&1\end{pmatrix}, \quad \det P = 2, \quad P^{-1} = \frac{1}{2}\begin{pmatrix}1&-1\\1&1\end{pmatrix}
\]

\textbf{4.} $D^n = \begin{pmatrix}1^n&0\\0&3^n\end{pmatrix} = \begin{pmatrix}1&0\\0&3^n\end{pmatrix}$, donc :
\[
A^n = PD^nP^{-1} = \begin{pmatrix}1&1\\-1&1\end{pmatrix}\begin{pmatrix}1&0\\0&3^n\end{pmatrix}\frac{1}{2}\begin{pmatrix}1&-1\\1&1\end{pmatrix}
\]
\[
= \frac{1}{2}\begin{pmatrix}1&3^n\\-1&3^n\end{pmatrix}\begin{pmatrix}1&-1\\1&1\end{pmatrix}
= \frac{1}{2}\begin{pmatrix}1+3^n & -1+3^n \\ -1+3^n & 1+3^n\end{pmatrix}
\]

\textbf{5.} Pour $n=5$ : $3^5 = 243$ :
\[
A^5 = \frac{1}{2}\begin{pmatrix}244 & 242 \\ 242 & 244\end{pmatrix} = \begin{pmatrix}122 & 121 \\ 121 & 122\end{pmatrix}
\]
\end{boitecorrige}

% ================================================================
\subsection*{Exercice 18 - Système linéaire compatible indéterminé \niveau{\moyen}}
% ================================================================
\newpage
\begin{boiteexercice}
Résoudre par la méthode de Gauss-Jordan le système suivant à 4 inconnues :
\[
\left\{\begin{aligned}
x_1 + 2x_2 - x_3 + x_4 &= 2 \\
2x_1 + 4x_2 + x_3 - x_4 &= 5 \\
x_1 + 2x_2 + 2x_3 - 2x_4 &= 3
\end{aligned}\right.
\]

\begin{enumerate}
  \item Écrire la matrice augmentée et effectuer les opérations élémentaires sur les lignes.
  \item Identifier le rang du système et les inconnues libres (paramètres).
  \item Exprimer les inconnues liées en fonction des paramètres.
  \item Écrire la solution générale sous forme vectorielle.
\end{enumerate}
\end{boiteexercice}

\begin{boitecorrige}
\textbf{1. Matrice augmentée et réduction :}
\[
\left(\begin{array}{cccc|c}1&2&-1&1&2\\2&4&1&-1&5\\1&2&2&-2&3\end{array}\right)
\]
$L_2 \leftarrow L_2 - 2L_1$ ; $L_3 \leftarrow L_3 - L_1$ :
\[
\left(\begin{array}{cccc|c}1&2&-1&1&2\\0&0&3&-3&1\\0&0&3&-3&1\end{array}\right)
\]
$L_3 \leftarrow L_3 - L_2$ :
\[
\left(\begin{array}{cccc|c}1&2&-1&1&2\\0&0&3&-3&1\\0&0&0&0&0\end{array}\right)
\]
$L_2 \leftarrow L_2/3$ ; $L_1 \leftarrow L_1 + L_2$ :
\[
\left(\begin{array}{cccc|c}1&2&0&0&7/3\\0&0&1&-1&1/3\\0&0&0&0&0\end{array}\right)
\]

\textbf{2.} Le rang est $r=2$. Les colonnes pivots sont 1 et 3. Les inconnues libres (paramètres) sont $x_2 = s$ et $x_4 = t$ ($s,t\in\mathbb{R}$).

\textbf{3.} Des équations réduites :
\begin{align*}
x_1 + 2s &= 7/3 &\Rightarrow& \quad x_1 = 7/3 - 2s \\
x_3 - t &= 1/3  &\Rightarrow& \quad x_3 = 1/3 + t
\end{align*}

\textbf{4. Solution générale sous forme vectorielle :}
\[
\begin{pmatrix}x_1\\x_2\\x_3\\x_4\end{pmatrix}
= \begin{pmatrix}7/3\\0\\1/3\\0\end{pmatrix}
+ s\begin{pmatrix}-2\\1\\0\\0\end{pmatrix}
+ t\begin{pmatrix}0\\0\\1\\1\end{pmatrix}, \quad s,t\in\mathbb{R}
\]

L'espace solution est un \textbf{sous-espace affine} de dimension 2 (plan affine dans $\mathbb{R}^4$). La solution particulière est $\vec{x}_0 = (7/3, 0, 1/3, 0)^T$ et la solution homogène est engendrée par les vecteurs $(-2,1,0,0)^T$ et $(0,0,1,1)^T$.
\end{boitecorrige}

\ifdefined\inMainDoc\else\end{document}\fi
